\begin{proposition}[Refuted: every admissible-length ordered necklace is orderable]
\label{C:ordered-necklace-short-warning}
The sufficiently-large-$k$ hypothesis in
Proposition~\ref{C:ordered-necklace} cannot be replaced, for arbitrary
increasing maps of finite linear orders, by the condition that $k$
bounds the sizes of their nontrivial point paths. The failure can
occur at $n=3$, $s=45$, $k=11$, satisfying the numerical C.1 condition
for $a_0=2$, $a_1=8/5$.
The pair below is not asserted to occur among the actual generated
arrows of C.1, so this does not refute the universal statement of
Part C.1 or the selected construction already proved.
\end{proposition}

\begin{proof}
For any integer $k\geq2$, put
\[
 d=2k,\qquad s=2d+1=4k+1,\qquad M=d(d+1).
\]
Give the finite sets
\[
 S=\{0,\ldots,M\},\qquad T=\{0,\ldots,2M+s\}
\]
their natural linear orders, and define
\[
 f,g:S\rightrightarrows T,
 \qquad f(r)=2r,\quad g(r)=2r+s.
\]
Both are increasing injections. Their images are disjoint, since
$s$ is odd, and their nontrivial point relation consists entirely
of two-vertex paths. Thus the repaired target with this $k$ and
actual path size two is defined and satisfies the usual neighboring
identities.

Write those identities as
$\alpha_i\lambda_i=\alpha_{i+1}\mu_i$, with cyclic indices and
\[
 (\lambda_i,\mu_i)=
 \begin{cases}
 (g,f),&0\leq i<d,\\
 (g,g),&i=d,\\
 (f,g),&d<i<2d,\\
 (f,f),&i=2d.
 \end{cases}
\]
We give an explicit cycle of forced strict inequalities in the
repaired target. Begin with $r_{-1}=M$ and define
\[
 r_i=
 \begin{cases}
 r_{i-1}-d,&0\leq i\leq d,\\
 r_{i-1}+d+1,&d<i\leq2d.
 \end{cases}
\]
All these indices lie in $S$: the first part decreases from $M-d$
to zero, and the second increases from $d+1$ to $M$.
In particular $r_{2d}=r_{-1}$, so this is a cyclic list.

For $0\leq i\leq d$, the incoming and outgoing old points in copy
$i$ are $f(r_{i-1})$ and $g(r_i)$, and
\[
 g(r_i)=2(r_{i-1}-d)+(2d+1)=f(r_{i-1})+1.
\]
For $d<i\leq2d$, they are $g(r_{i-1})$ and $f(r_i)$, and
\[
 f(r_i)=2(r_{i-1}+d+1)=g(r_{i-1})+1.
\]
Thus in every copy the incoming point is strictly smaller than the
outgoing point, in fact immediately precedes it in $T$. Equivalently,
\[
 \mu_{i-1}(r_{i-1})<\lambda_i(r_i)
 \qquad(0\leq i<s),
\]
where the lower junction uses $\mu_{2d}=f$.
If every $\alpha_i$ preserved the old linear order, these inequalities
and the neighboring equalities would give
\[
 \alpha_0\mu_{s-1}(r_{s-1})
 <\alpha_0\lambda_0(r_0)
 =\alpha_1\mu_0(r_0)
 <\cdots
 <\alpha_{s-1}\lambda_{s-1}(r_{s-1})
 =\alpha_0\mu_{s-1}(r_{s-1}).
\]
This is impossible in a linear order. No order on this repaired
target therefore makes all its copy maps increasing.

Taking $k=11$, $n=3$ gives $d=22$, $M=506$, and $s=45$.
The source has $507$ points, the old target has $1058$ points,
and the repaired target has $24795$ points. Its forced precedence
cycle has $45$ edges. The standalone program
\texttt{akr\_ordering\_counterexample.py} constructs the repaired
maps, verifies the source path sizes and all neighboring identities,
and checks each consecutive-point comparison and equality in that
cycle. The output and JSON certificate give all $45$ comparisons.
Since $45/(3\cdot3)=5=(1-(8/5)/2)^{-1}$, these are numerically
admissible C.1 parameters for a first-step ratio, though no realization
of this particular source pair in the prescribed construction is
claimed.
\end{proof}
